\documentclass[11pt,a4paper]{article}
\usepackage[T1]{fontenc}
\usepackage[utf8]{inputenc}
\usepackage{lmodern,amsmath,amssymb,mathtools,microtype}
\usepackage[margin=25mm]{geometry}
\usepackage{fancyhdr}
\usepackage[hidelinks]{hyperref}
\usepackage{xurl}
\setlength{\parskip}{0.42em}
\setlength{\parindent}{0pt}
\setlength{\emergencystretch}{2em}
\providecommand{\tightlist}{\setlength{\itemsep}{0pt}\setlength{\parskip}{0pt}}
\pagestyle{fancy}
\fancyhf{}
\fancyhead[L]{\small Four-symbol reversal detection}
\fancyhead[R]{\small No.39 / v0.4}
\fancyfoot[C]{\thepage}
\setlength{\headheight}{14pt}
\title{Four-symbol reversal detection\\on a homogeneous six-state cycle}
\author{Ryutaro Yonezu\\\small Yonezu Research Lab\\\small\href{mailto:yonezaemon1+research@gmail.com}{\texttt{yonezaemon1+research@gmail.com}}}
\date{25 October 2026 \quad Version 0.4}
\begin{document}
\maketitle
\begin{abstract}
We study a stationary continuous-time nearest-neighbor walk on a homogeneous
six-cycle, observed at equal time intervals through the binary labeling
$001011$. We give an exact formula for the probability difference between
the consecutive words $0011$ and $1100$, and prove that its remaining
polynomial factor is strictly positive. Thus this one length-four pair
detects every nonresonant time-reversal asymmetry in the specified model.
At the resonant intervals the entire observation process is reversible.
Length four is sharp because stationary binary words of length at most
three are reversal symmetric. Broader reconstruction implications are
established prior context; historical novelty of this explicit refinement
remains inconclusive.
\end{abstract}
\begin{quote}\small
Unrefereed, AI-assisted research preprint. The mathematical body follows the
audited v0.4 note. A received AI review reports independent proof derivation;
its reflection condition was subsequently marked PASS, and the remaining
FTY source condition was checked author-side against the original appendix.
This is not human peer review. Literature coverage is partial; no priority
or first-discovery claim is made. The final Howard and Matzinger--Rolles
full-text comparisons remain incomplete.
\end{quote}
\hypertarget{model-and-statement}{%
\section{Model and statement}\label{model-and-statement}}

Let \(X(t)\) be a continuous-time nearest-neighbor Markov chain on
\(\mathbb{Z}/6\mathbb{Z}\), with clockwise rate \(a>0\) and
counterclockwise rate \(b>0\) at every vertex. The initial distribution
is uniform. Set \(f=(0,0,1,0,1,1)\) in clockwise vertex order and
\(Y_k=f(X(k\tau))\), where \(\tau>0\). A word probability \(p(w)\)
refers to the stationary law of consecutive Y's, not the reliability of
a single observed word.

Put

\[r=e^{-(a+b)\tau/2},\qquad
\theta=\frac{\sqrt3}{2}(a-b)\tau,\qquad c=\cos\theta.\]

\textbf{Candidate theorem.} Define

\[G(r,c)=1-2r^2+4r^6+r^8+4r^2(1-r^2)c^2-2r^3(1+r^4)c.\]

Then

\[p(0011)-p(1100)=-\frac{\sqrt3}{54}r^3\sin\theta\,G(r,\cos\theta),\]

and \(G(r,c)>0\) for \(0<r<1\) and \(-1\le c\le1\). Consequently:

\begin{enumerate}
\def\labelenumi{\arabic{enumi}.}
\tightlist
\item
  If \(\sin\theta\) is nonzero, this one pair of length-four words
  witnesses time-reversal asymmetry. No length-one, two, or three
  consecutive words can witness such asymmetry for any stationary binary
  process.
\item
  If \(\sin\theta=0\), the entire sampled hidden chain, and therefore
  its binary observation process, is time-reversal symmetric at every
  word length.
\end{enumerate}

For \(a\ne b\) the blind spacings are exactly

\[\tau=\frac{2\pi m}{\sqrt3\,|a-b|},\qquad m=1,2,3,\ldots.\]

For \(a=b\) all spacings are reversible.

The candidate contribution is the explicit \textbf{length-four
sufficiency} and its positive prefactor. The all-length dichotomy also
follows from established scenery reconstruction, as explained in Section
6; it is not an independent novelty claim.

\hypertarget{transition-kernel-and-the-four-symbol-identity}{%
\section{Transition kernel and the four-symbol
identity}\label{transition-kernel-and-the-four-symbol-identity}}

For the row-generator convention,

\[ (Qh)(i)=a(h(i+1)-h(i))+b(h(i-1)-h(i)).\]

The characters \(\exp(2\pi\mathrm{i}ki/6)\) have eigenvalues

\[\mu_k=(a+b)(\cos(2\pi k/6)-1)+\mathrm{i}(a-b)\sin(2\pi k/6).\]

Thus the six eigenvalues of \(P=\exp(\tau Q)\), in character order, are

\[1,\ r(c+\mathrm{i}s),\ r^3(c+\mathrm{i}s),\ r^4,
\ r^3(c-\mathrm{i}s),\ r(c-\mathrm{i}s),\qquad s=\sin\theta.\]

Let \(K_j=P(0,j)\) and \(d=\sqrt{3}\). Fourier inversion gives

\[
\begin{aligned}
6K_0&=1+r^4+2rc+2r^3c,\\
6K_1&=1-r^4+rc+drs-r^3c+dr^3s,\\
6K_2&=1+r^4-rc+drs-r^3c-dr^3s,\\
6K_3&=1-r^4-2rc+2r^3c,\\
6K_4&=1+r^4-rc-drs-r^3c+dr^3s,\\
6K_5&=1-r^4+rc-drs-r^3c-dr^3s.
\end{aligned}
\]

For \(D_0=\operatorname{diag}(1-f)\) and \(D_1=\operatorname{diag}(f)\),

\[p(w)=\frac16\mathbf1^T D_{w_0}P D_{w_1}P\cdots
 P D_{w_{\ell-1}}\mathbf1.\]

Substitution for the two length-four words and use of \(c^2+s^2=1\)
gives the displayed probability difference. This step can be checked by
direct finite expansion: exact\_algebra.py enumerates microscopic paths,
works over \(\mathbb{Q}[r,c,s,d]/(d^2-3)\), and reduces the difference
modulo \(c^2+s^2-1\). The coefficient remainder is identically zero. No
numerical approximation enters this identity.

\hypertarget{strict-positivity-of-g}{%
\section{Strict positivity of G}\label{strict-positivity-of-g}}

View G as a quadratic in c.~Its leading coefficient is \(4r^2(1-r^2)>0\)
and its vertex is

\[c_* = \frac{r(1+r^4)}{4(1-r^2)}.\]

First suppose \(r^2\le3/4\) and put \(t=r^2\). The unrestricted
quadratic minimum is

\[G(r,c_*)=\frac{q(t)}{4(1-t)},\]

\[q(t)=4-12t+7t^2+16t^3-14t^4-4t^5-t^6.\]

On \(0\le t\le3/4\) set \(u=4t/3\). The degree-six Bernstein expansion
is

\[q(3u/4)=\sum_{j=0}^{6} B_j {6\choose j}u^j(1-u)^{6-j},\]

with

\[ (B_0,\ldots,B_6)=
\left(4,\frac52,\frac{101}{80},\frac58,
\frac{403}{640},\frac{443}{512},\frac{535}{4096}\right).\]

All seven coefficients are strictly positive and the Bernstein basis
sums to one. Hence \(q(t)>0\) throughout this interval, proving \(G>0\)
there for all real c.

Now suppose \(r^2\ge3/4\). The function

\[r(1+r^4)-4(1-r^2)\]

is strictly increasing for \(r>0\), and at \(r=\sqrt{3}/2\) has value
\(25\sqrt{3}/32-1>0\). Thus \(c_*>1\). On \(-1\le c\le1\), G decreases
as c increases, so it suffices to consider \(c=1\). Exact factorization
gives

\[G(r,1)=(1-r)^2
 (r^6+3r^4+6r^3+5r^2+2r+1)>0\quad(0<r<1).\]

This completes the positivity proof. The coefficient identity, vertex
minimum polynomial, and endpoint factorization are checked over
rationals by verify.py.

\hypertarget{sharp-word-length-and-the-resonant-case}{%
\section{Sharp word length and the resonant
case}\label{sharp-word-length-and-the-resonant-case}}

For any stationary binary process, the length-two identity
\(p(01)=p(10)\) follows from equality of one-site marginals. At length
three, stationarity of the word 00 gives \(p(001)=p(100)\), and
stationarity of 11 gives \(p(011)=p(110)\). All other length-three words
are palindromes. These facts prove the lower bound four without any
Markov assumption on Y. This lower bound is established prior art
{[}1{]}.

If \(\sin\theta=0\), all \(K_j=K_{-j}\), so the sampled chain has
symmetric P. For every hidden path \((i_0,\ldots,i_\ell)\), uniform
stationarity and \(P=P^T\) make its probability equal to that of its
reversed path. Summing over hidden paths with any specified binary
labels proves observed reversal symmetry at every length. If
\(\sin\theta\ne0\), Section 3 makes the displayed four-word difference
nonzero, establishing the converse.

\hypertarget{labeling-and-interpretation-limits}{%
\section{Labeling and interpretation
limits}\label{labeling-and-interpretation-limits}}

A binary labeling f is reflection-symmetric when \(f(i)=f(k-i)\) for
some k. For uniform stationarity, the reversed hidden transition matrix
is \(P_{a,b}^{T}=P_{b,a}\), because \(Q_{a,b}^{T}=Q_{b,a}\); hence the
time-reversed observation law at rates (a,b) is the forward observation
law at rates (b,a). The reflection \(\sigma(i)=k-i\) preserves the
uniform initial law, interchanges the two rates, and preserves the
observed label whenever \(f(i)=f(\sigma(i))\). Thus in this case the
forward laws at (a,b) and (b,a) agree, proving reversibility of the
observations at every \(\tau\), irrespective of drift.

Among the 64 labelings on six vertices, 52 have such a reflection and 12
do not. The latter form the two rotational orbits represented by 001011
and 001101. The second is a rotation of the reversed first labeling.
Rotation preserves word probabilities, and reflection interchanges a,b.
Thus the theorem covers every chiral labeling on the six-cycle, with a
possible sign reversal of the witness.

All binary labelings on cycles of lengths 3,4,5 have a reflection. These
small counts are exhaustively verified. The minimal size six is only
about \textbf{homogeneous nearest-neighbor cycles with deterministic
binary vertex labels}. It is not a minimal-state claim about general
hidden Markov models.

This note addresses stationary statistical time reversal of sampled
paths. It does not say that thermodynamic nonequilibrium is impossible
to infer from other structural information. In particular, at odd
resonances the sampled matrix has negative eigenvalues, obstructing a
reversible continuous-time embedding. At even resonances it equals the
sampled propagator of the symmetric ring with rates \(a'=b'=(a+b)/2\).

Nor does a length-four witness mean that four recorded bits suffice for
a reliable empirical decision. The probabilities need repeated
estimation. Signals may be extremely small near resonances or at large
\(\tau\). No finite-sample detection guarantee or dissipation estimator
is claimed.

\hypertarget{established-implications-and-remaining-novelty-question}{%
\section{Established implications and remaining novelty
question}\label{established-implications-and-remaining-novelty-question}}

For Finucane--Tamuz--Yaari {[}2{]}, the source fixed here is the
author-hosted, 17-page manuscript dated 30 March 2014, not an
unspecified arXiv version. Definition 1.1 (printed p.2) uses shifts
only. Theorem 1.2 and footnote 3 (printed p.3) distinguish this
convention from Matzinger--Lember's shift-and-reflection equivalence.
The explicit shift-only statement we use is Theorem Appendix A.1
(printed p.14), with its proof completed in Appendix A.3 (printed p.16).
It requires the Fourier coefficients to be pairwise distinct \textbf{as
complex numbers}. A reflection that is not a shift can therefore be
distinguished. The original Matzinger--Lember theorem is not silently
substituted for this shift-only statement.

Here \(0<r<1\). The four groups

\[\{1\},\quad\{re^{i\theta},re^{-i\theta}\},\quad
\{r^3e^{i\theta},r^3e^{-i\theta}\},\quad\{r^4\}\]

have distinct moduli \(1,r,r^3,r^4\). Within either conjugate pair, the
two complex numbers differ exactly when \(\sin\theta\ne0\). Thus all six
Fourier coefficients are pairwise distinct as complex numbers in the
nonresonant case; their six absolute values are not pairwise distinct.
Applying that theorem distinguishes the chiral labeling from its
reflection. At \(\sin\theta=0\) the direct symmetry argument applies.
The all-length dichotomy is consequently a known consequence, whereas
the explicit one-pair, length-four formula is the candidate refinement.

Periodic-sampling ambiguity itself is also known; Bauer--Seifert--van
der Meer {[}3{]}, Section 3.2.2, give different
equilibrium/nonequilibrium generators with one sampled propagator. The
general three-word obstruction is {[}1{]}. The remaining novelty
question is narrowly whether the explicit sharp four-consecutive-symbol
result on this minimal chiral homogeneous ring has already been proved,
or follows immediately from another known formula. See PRIOR\_ART.md for
the recovered Matzinger--Lember manuscript and its version/access
limits. The scoped search has not settled historical novelty. REFERENCES
are PARTIAL and NOVELTY is INCONCLUSIVE. An independent AI proof-review
report has been received within the scope above; human peer review has
not been performed. The direct proof in Sections 2-5 does not depend on
FTY. Further source-access details are in PRIOR\_COMPARISON.md.

\hypertarget{references}{%
\section*{References}\label{references}}

{[}1{]} L. Arola-Fernandez and L. Lacasa (2023), \emph{Irreversibility
of symbolic time series: a cautionary tale}. Phys. Rev.~E 108, 014201.
\url{https://arxiv.org/pdf/2303.11868}
\url{https://doi.org/10.1103/PhysRevE.108.014201}

{[}2{]} H. Finucane, O. Tamuz and Y. Yaari (2014), \emph{Scenery
reconstruction on finite abelian groups}. Stochastic Processes and their
Applications 124, 2754--2770. The source inspected is the author-hosted
manuscript dated 30 March 2014, 17 pages:
\url{https://tamuz.caltech.edu/papers/scenery.pdf}
\url{https://doi.org/10.1016/j.spa.2014.03.012}

{[}3{]} M. T. Bauer, U. Seifert and J. van der Meer (2025),
\emph{Stroboscopic measurements in Markov networks: exact generator
reconstruction vs. thermodynamic inference}. J. Phys. A 58, 125001.
\url{https://arxiv.org/html/2412.15642v2}
\url{https://doi.org/10.1088/1751-8121/adbc52}

\bigskip
\noindent\small Copyright \textcopyright\ 2026 Ryutaro Yonezu. All rights reserved.
No general reuse or redistribution license is granted.
\end{document}
